
\documentclass{article}
\usepackage{colortbl}
\usepackage{makecell}
\usepackage{multirow}
\usepackage{supertabular}

\begin{document}

\newcounter{utterance}

\centering \large Interaction Transcript for game `wordle', experiment `medium\_frequency\_words\_no\_clue\_no\_critic', episode 10 with NEW\_Reflect\_Agent.
\vspace{24pt}

{ \footnotesize  \setcounter{utterance}{1}
\setlength{\tabcolsep}{0pt}
\begin{supertabular}{c@{$\;$}|p{.15\linewidth}@{}p{.15\linewidth}p{.15\linewidth}p{.15\linewidth}p{.15\linewidth}p{.15\linewidth}}
    \# & \multicolumn{2}{c}{Player} && \multicolumn{2}{c}{Game Master} \\
    \hline

    \theutterance \stepcounter{utterance}  
    & & & \multicolumn{4}{p{0.6\linewidth}}{
        \cellcolor[rgb]{0.9,0.9,0.9}{
            \makecell[{{p{\linewidth}}}]{
                \texttt{\tiny{[P1$\langle$GM]}}
                \texttt{You are a language wizard who likes to guess words by using the given rules.} \\
\\ 
\texttt{Welcome to Wordle! You have six attempts to guess the target word, a valid English word of five lowercase letters (a{-}z). Please use the tags "explanation:" and "guess:" to provide a concise explanation for each guess.} \\
\\ 
\texttt{For instance, if your guess is "apple", your response should be} \\
\texttt{explanation: this is a common five{-}letter English word, and I am starting my guess with this word.} \\
\texttt{guess: apple} \\
\\ 
\texttt{After each guess, your answer will be validated, and you will receive feedback indicating which letters are correct (green), which letters are correct but in the wrong position (yellow), and which letters are incorrect (red). This feedback can be useful in determining which letters to include or exclude in your next guess.} \\
\\ 
\texttt{For example, the feedback for "apple" might be:} \\
\texttt{guess\_feedback: a<yellow> p<yellow> p<green> l<yellow> e<red>} \\
\\ 
\texttt{The explanation should contain details about how the guess\_feedback is used to arrive at a new guess.} \\
\\ 
\texttt{Do not generate any other text.} \\
\\ 
\texttt{Let's begin with your first guess.} \\
            }
        }
    }
     \\ \\

    \theutterance \stepcounter{utterance}  
    & \multicolumn{4}{p{0.45\linewidth}}{
        \cellcolor[rgb]{0.95,0.95,0.95}{
            \makecell[{{p{\linewidth}}}]{
                \texttt{\tiny{[P1$|internal]}}
                \texttt{Reflections from previous episodes:} \\
\texttt{{-}} \\
\texttt{ Failed. I repeated the same guess multiple times despite identical feedback, wasting attempts. The feedback showed only 'e' in position 2 was green, but I didn't adjust my guess to test other letters for remaining positions. Next time, I would use each guess to systematically test new letters based on feedback, not repeat the same word.} \\
\texttt{{-}} \\
\texttt{ I failed because I repeated the same guess multiple times despite receiving identical feedback. The feedback showed 'n' and 'e' were correct in positions 4 and 5, but I did not adjust my guess to incorporate new letters or positions. Next time, I would use the feedback to systematically eliminate letters and test new possibilities, ensuring each guess builds on previous information rather than stagnating.} \\
\texttt{{-}} \\
\texttt{ I failed to guess the target word. The first guess "crane" eliminated five common letters. The second guess "blind" gave three green letters (b, l, i) but left two incorrect letters (n, d). I should have used the confirmed "bli" and tested remaining positions with common vowels/consonants instead of guessing a full word. Next time, I would use a strategic guess to efficiently test possible letters for the last two slots.} \\
            }
        }
    }
    & & \\ \\

    \theutterance \stepcounter{utterance}  
    & \multicolumn{4}{p{0.45\linewidth}}{
        \cellcolor[rgb]{0.95,0.95,0.95}{
            \makecell[{{p{\linewidth}}}]{
                \texttt{\tiny{[P1$|internal]}}
                \texttt{{[}attempt 1 – valid{]}} \\
\texttt{explanation: I start with a common five{-}letter word that includes frequently used vowels and consonants to gather initial feedback about which letters are in the target word and their positions.} \\
\texttt{guess: crane} \\
            }
        }
    }
    & & \\ \\

    \theutterance \stepcounter{utterance}  
    & \multicolumn{4}{p{0.6\linewidth}}{
        \cellcolor[rgb]{0.561,0.737,0.561}{
            \makecell[{{p{\linewidth}}}]{
                \texttt{\tiny{[P1$\rangle$GM]}}
                \texttt{explanation: I start with a common five{-}letter word that includes frequently used vowels and consonants to gather initial feedback about which letters are in the target word and their positions.} \\
\texttt{guess: crane} \\
            }
        }
    }
    & & \\ \\

    \theutterance \stepcounter{utterance}  
    & & & \multicolumn{2}{p{0.3\linewidth}}{
        \cellcolor[rgb]{0.9,0.9,0.9}{
            \makecell[{{p{\linewidth}}}]{
                \texttt{\tiny{[GM$|$GM]}}
                \texttt{attempts = 1} \\
\texttt{target\_word = upper} \\
\texttt{guess = crane} \\
\texttt{guess\_feedback = c<red> r<yellow> a<red> n<red> e<yellow>} \\
            }
        }
    }
    & & \\ \\

    \theutterance \stepcounter{utterance}  
    & & & \multicolumn{4}{p{0.6\linewidth}}{
        \cellcolor[rgb]{0.9,0.9,0.9}{
            \makecell[{{p{\linewidth}}}]{
                \texttt{\tiny{[P1$\langle$GM]}}
                \texttt{guess\_feedback: c<red> r<yellow> a<red> n<red> e<yellow>} \\
\\ 
\texttt{Provide your response only in this format.} \\
\texttt{explanation: details} \\
\texttt{guess: word} \\
            }
        }
    }
     \\ \\

    \theutterance \stepcounter{utterance}  
    & \multicolumn{4}{p{0.45\linewidth}}{
        \cellcolor[rgb]{0.95,0.95,0.95}{
            \makecell[{{p{\linewidth}}}]{
                \texttt{\tiny{[P1$|internal]}}
                \texttt{{[}attempt 1 – valid{]}} \\
\texttt{explanation: The feedback shows that 'r' and 'e' are present in the target word but in the wrong positions, while 'c', 'a', and 'n' are not in the word. I will avoid using 'c', 'a', and 'n' in future guesses and focus on placing 'r' and 'e' in different positions.} \\
\texttt{guess: rater} \\
            }
        }
    }
    & & \\ \\

    \theutterance \stepcounter{utterance}  
    & \multicolumn{4}{p{0.6\linewidth}}{
        \cellcolor[rgb]{0.561,0.737,0.561}{
            \makecell[{{p{\linewidth}}}]{
                \texttt{\tiny{[P1$\rangle$GM]}}
                \texttt{explanation: The feedback shows that 'r' and 'e' are present in the target word but in the wrong positions, while 'c', 'a', and 'n' are not in the word. I will avoid using 'c', 'a', and 'n' in future guesses and focus on placing 'r' and 'e' in different positions.} \\
\texttt{guess: rater} \\
            }
        }
    }
    & & \\ \\

    \theutterance \stepcounter{utterance}  
    & & & \multicolumn{2}{p{0.3\linewidth}}{
        \cellcolor[rgb]{0.9,0.9,0.9}{
            \makecell[{{p{\linewidth}}}]{
                \texttt{\tiny{[GM$|$GM]}}
                \texttt{attempts = 2} \\
\texttt{target\_word = upper} \\
\texttt{guess = rater} \\
\texttt{guess\_feedback = r<red> a<red> t<red> e<green> r<green>} \\
            }
        }
    }
    & & \\ \\

    \theutterance \stepcounter{utterance}  
    & & & \multicolumn{4}{p{0.6\linewidth}}{
        \cellcolor[rgb]{0.9,0.9,0.9}{
            \makecell[{{p{\linewidth}}}]{
                \texttt{\tiny{[P1$\langle$GM]}}
                \texttt{guess\_feedback: r<red> a<red> t<red> e<green> r<green>} \\
\\ 
\texttt{Provide your response only in this format.} \\
\texttt{explanation: details} \\
\texttt{guess: word} \\
            }
        }
    }
     \\ \\

    \theutterance \stepcounter{utterance}  
    & \multicolumn{4}{p{0.45\linewidth}}{
        \cellcolor[rgb]{0.95,0.95,0.95}{
            \makecell[{{p{\linewidth}}}]{
                \texttt{\tiny{[P1$|internal]}}
                \texttt{{[}attempt 1 – valid{]}} \\
\texttt{explanation: The feedback shows that 'e' and 'r' are in the correct positions (green), while 'r', 'a', 't' are incorrect or misplaced. Since 'e' and 'r' are confirmed in the last two positions, and 'r' is red in the first position, it means 'r' is not in the first position. The second letter is red, so 'a' is not in the second position. The third letter is red, so 't' is not in the third position. The first letter is red, so 'c' is not in the first position. The fourth letter is yellow, so 'n' is in the word but not in the fourth position. The fifth letter is yellow, so 'e' is in the word but not in the fifth position. However, 'e' is green in the fifth position, so it is correct. The word ends with 'er'. The first letter is not 'c', 'r', 'a', 't', or 'n'. The second letter is not 'a'. The third letter is not 't'. The fourth letter is not 'n'. The fifth letter is 'e'. The word must be a five{-}letter word ending in 'er', with 'e' in the fifth position and 'r' in the fourth position. The first letter is not 'c', 'r', 'a', 't', or 'n'. The second letter is} \\
            }
        }
    }
    & & \\ \\

    \theutterance \stepcounter{utterance}  
    & \multicolumn{4}{p{0.6\linewidth}}{
        \cellcolor[rgb]{0.561,0.737,0.561}{
            \makecell[{{p{\linewidth}}}]{
                \texttt{\tiny{[P1$\rangle$GM]}}
                \texttt{explanation: The feedback shows that 'e' and 'r' are in the correct positions (green), while 'r', 'a', 't' are incorrect or misplaced. Since 'e' and 'r' are confirmed in the last two positions, and 'r' is red in the first position, it means 'r' is not in the first position. The second letter is red, so 'a' is not in the second position. The third letter is red, so 't' is not in the third position. The first letter is red, so 'c' is not in the first position. The fourth letter is yellow, so 'n' is in the word but not in the fourth position. The fifth letter is yellow, so 'e' is in the word but not in the fifth position. However, 'e' is green in the fifth position, so it is correct. The word ends with 'er'. The first letter is not 'c', 'r', 'a', 't', or 'n'. The second letter is not 'a'. The third letter is not 't'. The fourth letter is not 'n'. The fifth letter is 'e'. The word must be a five{-}letter word ending in 'er', with 'e' in the fifth position and 'r' in the fourth position. The first letter is not 'c', 'r', 'a', 't', or 'n'. The second letter is} \\
            }
        }
    }
    & & \\ \\

    \theutterance \stepcounter{utterance}  
    & & & \multicolumn{2}{p{0.3\linewidth}}{
        \cellcolor[rgb]{0.9,0.9,0.9}{
            \makecell[{{p{\linewidth}}}]{
                \texttt{\tiny{[GM$|$GM]}}
                \texttt{Error: The response should contain the 'guess:' keyword exactly once.} \\
            }
        }
    }
    & & \\ \\

    \theutterance \stepcounter{utterance}  
    & & & \multicolumn{2}{p{0.3\linewidth}}{
        \cellcolor[rgb]{0.9,0.9,0.9}{
            \makecell[{{p{\linewidth}}}]{
                \texttt{\tiny{[GM$|$GM]}}
                \texttt{game\_result = ABORT} \\
            }
        }
    }
    & & \\ \\

\end{supertabular}
}

\end{document}
